\documentclass{article}
\usepackage{amsmath}
\usepackage{amsfonts}
\usepackage{amssymb}
\usepackage{graphicx}

\makeatletter
\newcommand{\authorseparator}{}
\newcommand{\authorlist}{}
\newcommand{\coauthor}[1]{\g@addto@macro\authorlist{\authorseparator{}\renewcommand{\authorseparator}{, }#1}}
\newcommand{\presentingauthor}[1]{\g@addto@macro\authorlist{\authorseparator{}\renewcommand{\authorseparator}{, }{\underline{#1}}}}

\setlength{\parindent}{0pt}
\date{}
\title{Chemotaxis, collective dynamics, and diffusive relays in neutrophil swarming}

\presentingauthor{Roma Gaidarov} % presenting author (underlined) -- replace with your name
\coauthor{Ariel Amir}          % replace / extend the author list as needed

\makeatother

\author{\authorlist{}} % Do not modify this field. Use the coauthor/presenting author commands above

\begin{document}

\maketitle

Neutrophils, the most abundant immune cells, swarm collectively around
infections and injuries. Recent experiments show that swarms are organized not
by static gradients but by self-propagating waves of chemoattractant:
stimulated cells relay the signal to their neighbours, while a negative
feedback loop stops the relay and limits its range [1]. A single
cell, by contrast, makes net progress toward the source only over a narrow
window of wave speeds [2]. The classical Keller--Segel model misses this,
predicting net motion away from the source at all speeds.

We model this in two parts. First, we treat the relay as an excitable
reaction-diffusion system and compare several inhibition mechanisms that could
stop the wave and set the swarm size; the resulting fronts are pushed waves.
Second, we couple cell motility to the cue field through the cell polarization, an
internal degree of freedom that sets the direction of motion and relaxes on its
own timescale. This lag breaks the symmetry of a passing pulse and gives a net
drift, reproducing the experimental observations of [2]. Combining the two, where cells both
make the wave and move in it, gives a self-consistent picture of neutrophil
swarming.

%break and start a new line to separate the references from the main text
[1] E.~Strickland, D.~Pan, C.~Godfrey, J.~S.~Kim, A.~Hopke, W.~Ji, M.~Degrange,
B.~Villavicencio, M.~K.~Mansour, C.~S.~Zerbe, D.~Irimia, A.~Amir, and
O.~D.~Weiner, \textit{Self-extinguishing relay waves enable homeostatic control
of human neutrophil swarming}, Developmental Cell {\bf 59}, 2659, (2024).

%break and start a new line to separate the references from the main text
[2] M.~Ishida, M.~Uwamichi, A.~Nakajima, and S.~Sawai,
\textit{Traveling-wave chemotaxis of neutrophil-like HL-60 cells},
Molecular Biology of the Cell {\bf 36}, ar17, (2025).

%sample figure -- replace with a snapshot of the interactive tool, or remove
%\begin{figure}[h]
%\begin{center}
%\includegraphics[width = 0.5\columnwidth]{figure.png}
%\caption{Caption.\label{figure}}
%\end{center}
%\end{figure}

\end{document}
